%% --- Project Summary (1 page) ---
%% NSF requires three labeled subsections: Overview, Intellectual Merit, Broader Impacts.

\begin{center}
  {\large\bfseries Project Summary}\\[0.3em]
  {\large\bfseries \sysname: AI-Enabled Digital Twins for Continuous Red Teaming of Cyberphysical Systems}
\end{center}

\noindent\textbf{Overview.}
Cyberphysical systems (\CPSs) underpin the Nation's critical
infrastructure---power grids, water plants, hospitals, manufacturing.
Yet assessing their security posture remains a largely manual, ad hoc,
and infrequent activity.
\CPSs are inherently difficult to analyze: they are heterogeneous
mixtures of commodity computing, specialized controllers, and legacy
hardware; they are dynamic in both infrastructure and usage; and they
execute processes that must not be disrupted. A network scan that is
routine in an enterprise environment can, in a chemical plant, trigger a
safety-critical shutdown; a fuzzing campaign against a medical device
can injure a patient. As a result, security assessments are scheduled
rarely and performed by scarce human experts, while adversaries---now
increasingly augmented with AI---operate continuously.

This project proposes \sysname, an agentic framework that enables the
\emph{continuous, in-depth, and non-disruptive} security assessment of
cyberphysical systems. \sysname uses cooperating teams of AI agents to
(i)~ingest documentation, configuration, and runtime signals from a
target \cps into a structured and vectorized knowledge base;
(ii)~passively observe the system to build an evolving model of its
components, software, and flows; (iii)~synthesize a \emph{digital twin}
that is a minimally representative yet faithful replica of the original;
(iv)~continuously red-team the twin using component-level analyses
(fuzzing, static analysis, symbolic execution) composed into multi-step
attack plans; and (v)~surface prioritized findings to human analysts,
who validate them on the original system and whose feedback refines
the twin. Four agent roles---Knowledge Base, Observer, Replication,
and Red Team---operate as a feedback loop rather than a linear
pipeline, so that twin and analysis co-evolve with the real system.

\smallskip
\noindent\textbf{Intellectual Merit.}
\sysname advances the state of the art in three ways. First, the project
develops new techniques for automatically constructing a living model of
a \cps from heterogeneous, partial, and sometimes inconsistent evidence,
combining graph-structured knowledge with natural-language
representations to support agent reasoning. Second, it formulates digital
twin synthesis as a \emph{controlled-fidelity replication problem}, in
which diversity of component classes, topological patterns, and data
flows are explicit quantifiable parameters; this makes it feasible to
replicate large systems at a useful level of abstraction rather than
requiring an infeasible bit-for-bit copy. Third, the project builds on
the PIs' prior work on firmware re-hosting (Conware, HALucinator), AI
planning for privilege escalation (ChainReactor), and autonomous
vulnerability discovery (Artiphishell, developed for the DARPA AIxCC
competition) to compose single-component vulnerabilities into end-to-end
multi-step attacks whose plans are expressed in a structured planning
language and executed against the twin without risk to the real system.
Together, these contributions create a new scientific foundation for
continuous, automated, adversarial assessment of cyberphysical systems.

\smallskip
\noindent\textbf{Broader Impacts.}
The security of \CPSs is a matter of direct national importance:
attacks on pipelines, hospitals, water utilities, and industrial control
networks have already caused measurable harm to U.S.\ communities, and
the gap between attacker capability and defender capacity continues to
widen. By operating continuously and without disrupting operations,
\sysname can substantially raise the baseline resilience of critical
infrastructure sectors that cannot currently afford traditional
red-team engagements. The project will release its framework as open
source, following the PIs' established practice with widely-used
artifacts such as \texttt{angr}. The PIs will integrate the research
into UCSB graduate and undergraduate curricula, involve high-school,
REU, and ERSP students (including students from underrepresented
groups), and organize an annual \cps-focused CTF competition open to
the public. Industry collaboration, internships, and publication in
top-tier security venues will ensure that results transfer rapidly to
practitioners defending real systems.

\smallskip
\noindent\textbf{Keywords:} cyberphysical systems; digital twins;
AI agents; red teaming; critical infrastructure security.
